% Blueprint: numina_test
% Created by Numina

\begin{document}

\begin{definition}[Laplacian matrix]
\label{def:laplacian}
\lean{NuminafuseTest.Blueprint.laplacian}
\leanfile{NuminafuseTest/Blueprint.lean}
\leanok
Let \(G=(V,E)\) be a simple undirected graph on \(n\) vertices.
Let \(A(G)\in\mathbb{R}^{n\times n}\) be its adjacency matrix and let
\(D(G)\in\mathbb{R}^{n\times n}\) be the diagonal degree matrix, whose
\(i\)-th diagonal entry equals \(\deg_G(i)\).
The \emph{Laplacian matrix} of \(G\) is
\[
L(G)=D(G)-A(G).
\]
\end{definition}

\begin{definition}[Algebraic connectivity]
\label{def:algebraic_connectivity}
\lean{NuminafuseTest.Blueprint.algebraicConnectivity}
\leanfile{NuminafuseTest/Blueprint.lean}
\uses{def:laplacian}
The matrix \(L(G)\) is real symmetric and positive semidefinite, so its
eigenvalues are real and nonnegative.
We order them as
\[
0=\lambda_1(G)\le \lambda_2(G)\le \cdots \le \lambda_n(G).
\]
The second smallest eigenvalue \(\lambda_2(G)\) is called the
\emph{algebraic connectivity} of \(G\).
\end{definition}

\begin{definition}[Complete bipartite graph]
\label{def:complete_bipartite}
\lean{NuminafuseTest.Blueprint.completeBipartite}
\leanfile{NuminafuseTest/Blueprint.lean}
\leanok
Let \(a,b\ge 1\).
The \emph{complete bipartite graph} \(K_{a,b}\) is the simple undirected
graph whose vertex set can be partitioned as
\[
V = V_1 \sqcup V_2, \quad |V_1|=a,\; |V_2|=b,
\]
and whose edge set consists of all pairs \(\{u,v\}\) with \(u\in V_1\)
and \(v\in V_2\); there are no edges within \(V_1\) or within \(V_2\).
\end{definition}

\begin{lemma}[Algebraic connectivity of \(K_{2,n-2}\)]
\label{lem:lambda2_k2_nm2}
\lean{NuminafuseTest.Blueprint.lambda2_K2_nm2}
\leanfile{NuminafuseTest/Blueprint.lean}
\uses{def:algebraic_connectivity, def:complete_bipartite}
For every integer \(n\ge 4\),
\[
\lambda_2(K_{2,n-2}) = 2.
\]
\end{lemma}
\begin{proof}
Let
\[
K_{2,n-2}
\]
have bipartition
\[
V = V_1 \sqcup V_2,
\qquad
V_1=\{u_1,u_2\},
\qquad
V_2=\{v_1,\dots,v_{n-2}\}.
\]
Set
\[
m=n-2.
\]
Since \(n\ge 4\), we have \(m\ge 2\). Thus the graph is \(K_{2,m}\).

By the definition of the complete bipartite graph, every vertex in
\(V_1\) is adjacent to every vertex in \(V_2\), and there are no edges
inside \(V_1\) or inside \(V_2\). Hence
\[
\deg(u_1)=\deg(u_2)=m
\]
and
\[
\deg(v_i)=2
\qquad
\text{for every } i=1,\dots,m.
\]

Using the vertex order
\[
u_1,u_2,v_1,\dots,v_m,
\]
the Laplacian matrix \(L=L(K_{2,m})\) has the block form
\[
L=
\begin{pmatrix}
mI_2 & -J_{2,m}\\
-J_{m,2} & 2I_m
\end{pmatrix},
\]
where \(J_{r,s}\) denotes the \(r\times s\) all-ones matrix.
This follows from the definition \(L=D-A\): the diagonal entries are the
degrees, and the off-diagonal entry \(L_{ij}\) is \(-1\) exactly when the
corresponding vertices are adjacent.

We now compute the eigenvalues of \(L\).

First, since every row of a Laplacian matrix sums to zero, the all-ones
vector
\[
\mathbf{1}=(1,1,1,\dots,1)^T
\]
satisfies
\[
L\mathbf{1}=0.
\]
Thus \(0\) is an eigenvalue.

Next, consider a vector supported only on the second part \(V_2\), say
\[
x=(0,0,y_1,\dots,y_m)^T,
\]
where
\[
\sum_{i=1}^m y_i=0.
\]
For the first two coordinates, corresponding to \(u_1\) and \(u_2\), we
have
\[
(Lx)_{u_j}
=
m\cdot 0-\sum_{i=1}^m y_i
=
0
\qquad
(j=1,2),
\]
because the \(y_i\)'s sum to zero. For a vertex \(v_i\in V_2\), the only
neighbors are \(u_1\) and \(u_2\), whose coordinates in \(x\) are both
zero. Hence
\[
(Lx)_{v_i}
=
2y_i-0-0
=
2y_i.
\]
Therefore
\[
Lx=2x.
\]
The space of vectors \((y_1,\dots,y_m)\in\mathbb{R}^m\) satisfying
\(\sum_i y_i=0\) has dimension \(m-1\). Hence \(2\) is an eigenvalue of
\(L\) with multiplicity at least \(m-1=n-3\).

Similarly, consider the vector
\[
z=(1,-1,0,\dots,0)^T.
\]
For the first coordinate,
\[
(Lz)_{u_1}
=
m\cdot 1-\sum_{i=1}^m 0
=
m,
\]
and for the second coordinate,
\[
(Lz)_{u_2}
=
m\cdot (-1)-\sum_{i=1}^m 0
=
-m.
\]
For each \(v_i\in V_2\), we get
\[
(Lz)_{v_i}
=
2\cdot 0-1-(-1)
=
0.
\]
Thus
\[
Lz=mz.
\]
So \(m=n-2\) is also an eigenvalue.

It remains to find the eigenvalue on the two-dimensional subspace of
vectors that are constant on each bipartition class. Let
\[
x=(a,a,b,\dots,b)^T.
\]
For a vertex \(u_j\in V_1\),
\[
(Lx)_{u_j}
=
ma-\sum_{i=1}^m b
=
m(a-b).
\]
For a vertex \(v_i\in V_2\),
\[
(Lx)_{v_i}
=
2b-a-a
=
2(b-a).
\]
Therefore
\[
Lx
=
\bigl(m(a-b),m(a-b),2(b-a),\dots,2(b-a)\bigr)^T.
\]

We already know that when \(a=b\), this gives the eigenvector
\(\mathbf{1}\) with eigenvalue \(0\). To find the other eigenvalue in
this two-dimensional subspace, take \(x\) orthogonal to \(\mathbf{1}\).
That is, require
\[
2a+mb=0.
\]
One convenient choice is
\[
a=m,
\qquad
b=-2.
\]
Then
\[
x=(m,m,-2,\dots,-2)^T.
\]
For \(u_j\in V_1\),
\[
(Lx)_{u_j}
=
m(m-(-2))
=
m(m+2),
\]
while
\[
(m+2)x_{u_j}
=
(m+2)m
=
m(m+2).
\]
For \(v_i\in V_2\),
\[
(Lx)_{v_i}
=
2((-2)-m)
=
-2(m+2),
\]
while
\[
(m+2)x_{v_i}
=
(m+2)(-2)
=
-2(m+2).
\]
Hence
\[
Lx=(m+2)x.
\]
Since \(m=n-2\), this eigenvalue is
\[
m+2=n.
\]

We have now found the following eigenvalues:
\[
0,\qquad 2 \text{ with multiplicity } m-1,\qquad m,\qquad m+2.
\]
The dimensions add up to
\[
1+(m-1)+1+1=m+2=n,
\]
so these are all eigenvalues of \(L(K_{2,m})\). Replacing \(m\) by
\(n-2\), the Laplacian spectrum is
\[
0,\qquad 2^{(n-3)},\qquad n-2,\qquad n.
\]

Since \(n\ge 4\), we have
\[
n-2\ge 2.
\]
Therefore the second smallest Laplacian eigenvalue is
\[
\lambda_2(K_{2,n-2})=2.
\]
By the definition of algebraic connectivity, this proves that
\[
\lambda_2(K_{2,n-2})=2.
\]
\end{proof}

\begin{theorem}[Edge-count bound on algebraic connectivity]
\label{thm:lambda2_bound_m_2nm2}
\lean{NuminafuseTest.Blueprint.lambda2_bound_m_2nm2}
\leanfile{NuminafuseTest/Blueprint.lean}
\uses{def:algebraic_connectivity}
Let \(n\ge 4\), and let \(G\) be a simple undirected graph on \(n\)
vertices with exactly
\[
|E(G)|=2(n-2)
\]
edges.
Then
\[
\lambda_2(G)\le 2.
\]
\end{theorem}

\begin{theorem}[\(K_{2,n-2}\) maximizes algebraic connectivity for \(m=2(n-2)\)]
\label{thm:k2nminus2_maximizer}
\lean{NuminafuseTest.Blueprint.k2nminus2_maximizer}
\leanfile{NuminafuseTest/Blueprint.lean}
\uses{def:algebraic_connectivity, def:complete_bipartite, lem:lambda2_k2_nm2, thm:lambda2_bound_m_2nm2}
\leanok
Let \(n\ge 4\).
By Theorem~\ref{thm:lambda2_bound_m_2nm2}, the complete
bipartite graph \(K_{2,n-2}\) is a maximizer of algebraic connectivity
among all simple undirected graphs on \(n\) vertices with exactly
\(2(n-2)\) edges:
\[
\lambda_2(K_{2,n-2}) = 2,
\qquad
\lambda_2(G) \le \lambda_2(K_{2,n-2})
\quad
\text{for every such } G.
\]
\end{theorem}

\end{document}
